% MAS1004 Assignment 2: long report.
%
% You fill this in with your agent, in your repository. AGENTS.md tells the
% agent how. Red text in the PDF is what still has to be replaced.
%
% - Pictures go in figures/. Copy them from results/, which is not in git.
% - Outputs go in outputs/: copy experiments.csv, LOG.md and the outputs of
%   check_submission.py and curate.py count here. \outputfile{outputs/...}
%   prints a file exactly as it is. Every number in this report has to come
%   from such an output.
%
% To make the PDF: compress this folder into a zip, and on Overleaf choose
% New Project, Upload Project. Set Menu, Compiler to XeLaTeX, and Recompile.

\documentclass[11pt]{article}
\usepackage{mas1004}

\reporttitle{Assignment 2: Long report}
\studentname{Nordin Geels}
\studentid{G20260447}

\begin{document}
\makeheader

Leaderboard nickname: \textbf{Stroopwafel-Seoul}. Last releases:
\texttt{track-a-4} (Track A, recipe \texttt{ema\_no\_smooth}) and
\texttt{track-b-3} (Track B, \texttt{b\_mnv3\_512\_e8\_bal}).
\fillin{update if a later release is made}

All validation scores below are the balanced accuracy on
\texttt{data/starter/val} printed by \texttt{src/official.py} (Track A) and
\texttt{src/train\_b.py} (Track B), copied from \texttt{experiments.csv}.
The quotes are my own answers from \texttt{LOG.md}.

% ---------------------------------------------------------------------------
\section*{Track A}

\begin{center}\footnotesize
\begin{tabular}{lp{0.6\linewidth}r}
\toprule
tag & what changed from the run before & validation \\
\midrule
baseline & starter recipe, all 44,257 training photos, seed 0 / seed 1 & 0.496 / 0.495 \\
manual\_drop & 3,019 photos left out using my own control list (selection of 41,238) & 0.492 \\
no\_flip & no random left-right flip & 0.491 \\
less\_zoom & random crop scale 0.7--1.0 instead of 0.35--1.0 (built on no\_flip) & 0.478 \\
balanced & crop back to 0.35--1.0; every city drawn equally often (built on no\_flip) & 0.498 \\
ema & hand in a running average of the weights (EMA, decay 0.999) & 0.500 \\
ema\_no\_smooth & no label smoothing (was 0.1) & \textbf{0.507} \\
\bottomrule
\end{tabular}
\end{center}

\textbf{How much is chance?} The starter recipe trained twice with only a
different seed gave 0.496 and 0.495. ``The difference of only 0.001 shows that
the results barely change between the two random seeds. So the baseline seems
stable.'' Still, most of my Track A changes moved the score by less than 0.01,
so I treat those differences carefully.

\textbf{Removing photos (manual\_drop).} I expected a better score after taking
out photos ``that had nothing characteristic of the city'', but it went from
0.496 to 0.492: ``this is a bit of a shock''. See the next section.

\textbf{No flip.} Mirroring turns text around and puts traffic on the other
side of the road (Japan drives on the left, Korea and Taiwan on the right).
The score stayed the same (0.491): ``The model apparently relies little on such
small details and recognises a city mainly by the overall picture.''

\textbf{Less zoom} made it worse (0.478). ``With less random zooming the model
kept seeing roughly the same images'' and probably learned the training photos
themselves instead of the general features of a city.

\textbf{Balanced cities.} Busan (4,179) and Osaka (6,731) had far fewer
selected photos than Tokyo (13,055) and Seoul (10,834). Drawing every city
equally often raised Busan from 0.360 to 0.415 and Osaka from 0.344 to 0.427,
but lowered Seoul (0.597 $\to$ 0.552) and Tokyo (0.627 $\to$ 0.554); the total
went from 0.491 to 0.498. ``A fairer distribution does not automatically mean
that all cities are recognised better.''

\textbf{EMA} hardly changed anything (0.498 $\to$ 0.500): ``we do not know for
sure whether it came from EMA or from normal fluctuations.''

\textbf{No label smoothing} gave my best Track A score, 0.507 (0.500 before).
``Label smoothing was not really necessary in our case and possibly even
limited performance a little'', but the difference is small.

\textbf{Final recipe and selection:} \texttt{track\_a/recipe.py} as in
\texttt{ema\_no\_smooth} (balanced sampling, no flip, EMA, no label smoothing)
with my selection of 41,238 photos. Validation 0.507.
Leaderboard: \fillin{leaderboard score of track-a-4 when it appears}.
My first release (the starter recipe, validation 0.496) scored 0.403 on the
leaderboard.

% ---------------------------------------------------------------------------
\section*{Decisions I made from photos}

\picturefile{figures/mistakes_baseline_seed0_busan_00.png}

\fillin{What you looked at (the mistakes sheets of the baseline and the
look/manual sheets per city), what you saw in common, and the rule you wrote
before acting. In your own words.}

\picturefile{figures/manual_busan_00.png}

\textbf{What I changed.} I wrote a control list of 73 rules
(\texttt{track\_a/manual\_list.csv}): for each group of photos either drop all
of them or keep at most a few from a long series. The script
\texttt{src/manual\_list.py} turned it into \texttt{track\_a/selection.csv}
and \texttt{track\_a/dropped.csv}. Output of \texttt{curate.py count}:

\outputfile{outputs/curate_count.txt}

\textbf{Score before and after:} 0.496 / 0.495 (baseline, two seeds)
$\to$ 0.492 (manual\_drop). Leaving out 3,019 photos did not help; the
difference is within what I would expect from chance.

\textbf{My requests to the agent.} \fillin{Check and correct: I asked the
agent to show me contact sheets of each city and of the model's mistakes; I
gave it my own list of photo groups to drop (by city and Wikimedia ID range,
with a reason for each); I asked it to drop the three Busan videos with text
overlays; and I asked it to train the result as one measured run against the
baseline.}

% ---------------------------------------------------------------------------
\section*{Track B}

\begin{center}\footnotesize
\begin{tabular}{lp{0.6\linewidth}r}
\toprule
tag & what changed from the run before & validation \\
\midrule
b\_starter & \texttt{train\_b.py} as given: ResNet18, 224 px, 4 epochs & 0.489 \\
b\_mnv3\_448 & MobileNetV3-Large (ImageNet V2 weights) at 448 px & 0.562 \\
b\_mnv3\_448\_e6\_bal & 6 epochs instead of 4, every city drawn equally often & 0.574 \\
b\_mnv3\_512\_e8\_bal & 512 px (the maximum) and 8 epochs & \textbf{0.586} \\
\bottomrule
\end{tabular}
\end{center}

\textbf{Starting point.} ``It was useful to run the original code first,
because that way we had a clear starting point.'' It scored 0.489, close to
Track A's 0.496.

\textbf{Smaller model, bigger photos.} My idea was that ``bigger photos with
better quality'' let the model read text. MobileNetV3-Large at 448 px uses 0.85
billion multiply-adds, less than ResNet18 at 224 (1.8 billion), so it can look
at a photo four times as large. The score jumped from 0.489 to 0.562, the
largest step in this project. Because the model and the size changed together,
``we just cannot say yet which of the two changes had the most influence.''

\textbf{Longer and balanced.} 6 epochs with balanced cities: 0.574. Busan went
from 0.462 to 0.511 and Osaka from 0.436 to 0.519: ``not only did the total
score improve, but the differences between the cities also became smaller.''

\textbf{Largest size, longer.} 512 px and 8 epochs: 0.586, with Taipei at
0.700. 1.11 billion multiply-adds, within the limit of 2.0.

\textbf{Extra photos.} I did not collect extra photos; all Track B models were
trained on \texttt{data/starter/train} only.

\textbf{Validation next to leaderboard.} \fillin{leaderboard score of
track-b-3 when it appears}. Validation scores on the leaderboard are generally
much lower than on our validation photos (my starter recipe: 0.496 vs 0.403).
\fillin{Why you think they disagree, in your own words.}

\subsection*{My submission, as the grader sees it}
\outputfile{outputs/check_submission.txt}

% ---------------------------------------------------------------------------
\section*{What made the difference}

\fillin{In your own words. The numbers: in Track B, switching to MobileNetV3
at 448 px helped most (+0.073); balancing and longer training added +0.012,
and 512 px with 8 epochs another +0.012. In Track A no change moved the score
by more than 0.02: the best was ema\_no\_smooth (+0.011 over the baseline),
the worst less\_zoom ($-$0.013 compared with no\_flip).}

\section*{What I would do next with one more week}

``For Track A I would first try to distribute the 250,000 photos better over
the cities via \texttt{selection.csv}. [...] In \texttt{recipe.py} I would test
different learning rates, for example a lower learning rate [...]. I would
also try different photo edits, such as adjusting brightness and contrast.
For Track B I would train the model longer, for example 10 or 12 rounds
instead of 8. [...] I would test the image size, the number of training
rounds and the balancing of the cities separately.'' (full text at the end of
the lab notebook)

% ---------------------------------------------------------------------------
\section*{All runs}
\outputfile{outputs/experiments.csv}

\section*{Lab notebook}
\outputfile{outputs/LOG.md}

\end{document}
